% Who does what with ice2-data. Sits at the top of how-to/index.md, so a reader
% can find their own row before reading thirteen guide titles.
% Dashed ovals write; solid ovals only read.
\begin{tikzpicture}[
  node distance=5mm and 22mm,
  % Uniform width, so each actor's column reads as a column rather than as a
  % ragged pile. Without it the ellipses size to their labels and the left
  % edges wander by a centimetre.
  usecase/.append style={text width=34mm, minimum width=0mm},
  writecase/.append style={text width=34mm, minimum width=0mm},
]

% ---- actors, one per group of the guide index ---------------------------
\node[actor] (user) {Data user};
\node[actor, below=32mm of user] (author) {Library author};
\node[actor, below=46mm of author] (maint) {Catalogue\\maintainer};

% ---- what a data user does ----------------------------------------------
\node[usecase, right=of user, yshift=16mm]  (fetch)  {Fetch a collection};
\node[usecase, right=of user, yshift=4mm]   (cache)  {Put the cache somewhere};
\node[usecase, right=of user, yshift=-8mm]  (local)  {Use data already on disk};
\node[usecase, right=of user, yshift=-20mm] (check)  {Check and repair the cache};

\foreach \t in {fetch,cache,local,check} \draw[link] (user) -- (\t);

% ---- what a library author does -----------------------------------------
\node[usecase, right=of author, yshift=8mm]   (coll)  {Write a collections file};
\node[usecase, right=of author, yshift=-4mm]  (stage) {Stage uncatalogued data};
\node[usecase, right=of author, yshift=-16mm] (ci)    {Run it in CI};

\foreach \t in {coll,stage,ci} \draw[link] (author) -- (\t);

% ---- what a catalogue maintainer does -----------------------------------
\node[writecase, right=of maint, yshift=22mm]  (describe) {Describe a dataset};
\node[writecase, right=of maint, yshift=10mm]  (upload)   {Upload a dataset};
\node[writecase, right=of maint, yshift=-2mm]  (publish)  {Publish the catalogue};
\node[writecase, right=of maint, yshift=-14mm] (withdraw) {Withdraw a dataset};
\node[writecase, right=of maint, yshift=-26mm] (boot)     {Bootstrap a catalogue};

\foreach \t in {describe,upload,publish,withdraw,boot} \draw[link] (maint) -- (\t);

% The one relation worth drawing: a library can only name a dataset the
% catalogue already describes. Routed around the right-hand side -- straight
% through the middle would cross three unrelated use cases.
\draw[uses] (coll.east) -- ++(10mm,0) |- (describe.east)
  node[edgelabel, pos=0.28, right, xshift=1mm]{needs};

% ---- system boundary -----------------------------------------------------
\begin{scope}[on background layer]
  \node[boundarybox, fit=(fetch)(cache)(local)(check)(coll)(stage)(ci)
                        (describe)(upload)(publish)(withdraw)(boot)] (sys) {};
\end{scope}
\node[boundarylabel, anchor=north west] at (sys.north west) {ice2-data};

% ---- legend --------------------------------------------------------------
\node[usecase, text width=16mm, minimum height=7mm, font=\sffamily\scriptsize,
      anchor=north west, yshift=-9mm] (l1) at (sys.south west) {reads};
\node[writecase, text width=16mm, minimum height=7mm, font=\sffamily\scriptsize,
      right=5mm of l1] (l2) {writes};
\node[note, right=5mm of l2, align=left, text width=52mm]
  {every writing command lives under\\\texttt{ice2-data catalog ...}};

\end{tikzpicture}
